\documentclass[a4paper, 11pt]{article}

\usepackage[utf8]{inputenc}
\usepackage[spanish, activeacute, english]{babel}
\usepackage[left = 35mm, right = 25mm, top = 25mm, height = 206mm]{geometry}
\usepackage{fancyhdr}
\usepackage{graphicx}
\usepackage{lipsum}
\usepackage{natbib}


%----------------------------------------------------------------------------------------
%	ENCABEZADOS Y PIE DE PAGINA
%----------------------------------------------------------------------------------------

\pagestyle{fancy}
	\fancypagestyle{cuerpo}{
      \fancyfoot{}
      \fancyfoot[L]{Nombre Apellidos} %Autor
      \fancyfoot[C]{Asignatura} %Asignatura
      \fancyfoot[R]{Página: \thepage}
      \fancyhead{}
      \fancyhead[L]{Nombre del documento \\} %Nombre del documento
      \fancyhead[R]{\includegraphics[scale = 0.5]{logo_uco.jpg}} %Logo Universidad
      \renewcommand{\headrulewidth}{0.4pt}
      \renewcommand{\footrulewidth}{0.4pt}
      \setlength{\headheight}{65pt}
    }
    \fancypagestyle{indice}{
      \fancyfoot{}
      \fancyfoot[L]{Autor} %Autor
      \fancyfoot[R]{Asignatura} %Asignatura
      \fancyhead{}
      \fancyhead[L]{Nombre del documento \\} %Nombre del documento
      \fancyhead[R]{\includegraphics[scale = 0.5]{logo_uco.jpg}} %Logo
      \renewcommand{\headrulewidth}{0.4pt}
      \renewcommand{\footrulewidth}{0.4pt}
      \setlength{\headheight}{65pt}
    }

\newcommand{\excluirpagina}[1]{\addtocounter{page}{-1}} %Comando restar página.

\begin{document}
\selectlanguage{spanish}

%----------------------------------------------------------------------------------------
%	PAGINA DE TITULO
%----------------------------------------------------------------------------------------

\begin{titlepage}

\newcommand{\HRule}{\rule{\linewidth}{0.5mm}}

\center 
 
%----------------------------------------------------------------------------------------
%	ENCABEZADO DE LA PORTADA
%----------------------------------------------------------------------------------------

\textsc{\LARGE Universidad de Córdoba}\\[1.5cm] % Nombre de la Universidad
\textsc{\Large Escuela Técnica Superior de Ingeniería Agronómica y de Montes}\\[0.5cm] % Nombre de la escuela
\textsc{\large Trabajo Fin de Grado}\\[0.5cm] % Asignatura

%----------------------------------------------------------------------------------------
%	TITULOS
%----------------------------------------------------------------------------------------

\HRule \\[0.4cm]
{ \LARGE \bfseries Título del Trabajo Fin de Grado.}\\[0.4cm] % Titulo del Documento
\HRule \\[1.5cm]
 
%----------------------------------------------------------------------------------------
%	AUTOR
%----------------------------------------------------------------------------------------

\begin{minipage}{0.5\textwidth}
\begin{center}
\emph{Autores:}\\
Nombre Apellidos... % Nombre de Autor/es
\end{center}
\end{minipage}\\[1cm]
~
\begin{minipage}{0.5\textwidth}
\begin{center}
\emph{Tutor:}\\
Nombre Apellidos... % Nombre de Tutor/es
\end{center}
\end{minipage}\\[1cm]


%----------------------------------------------------------------------------------------
%	LOGOS
%----------------------------------------------------------------------------------------
\begin{minipage}{7cm}
\begin{flushleft}
\includegraphics[width=7.5cm]{logo_uco.jpg}
\end{flushleft}
\end{minipage}
~
\begin{minipage}{7cm}
\begin{flushright}
\includegraphics[width=7.5cm]{logo_uco.jpg}
\end{flushright}
\end{minipage}\\[3cm]

%----------------------------------------------------------------------------------------
%	FECHA
%----------------------------------------------------------------------------------------

{\large 10 de enero de 2017}\\[3cm] % Cambiar por \today para fecha actual


 
%----------------------------------------------------------------------------------------

\vfill 

\end{titlepage}

%----------------------------------------------------------------------------------------
%	RESUMENES Y ABSTRACT
%----------------------------------------------------------------------------------------

\newpage
\pagestyle{indice}
\thispagestyle{indice}
\excluirpagina

\begin{minipage}{\textwidth}
\begin{center}
\selectlanguage{english}
	\begin{abstract}
		Abstract in English %Resumen en Ingles
	\end{abstract}
\end{center}
\end{minipage}\\[90mm]
~
\begin{minipage}{\textwidth}
\begin{center}
\selectlanguage{spanish}
	\begin{abstract}
		Resumen en español %Resumen en Español
	\end{abstract}
\end{center}
\end{minipage}\\

%----------------------------------------------------------------------------------------
%	INDICES DE CONTENIDO 
%----------------------------------------------------------------------------------------

\newpage
\thispagestyle{fancy} 
\excluirpagina

\tableofcontents
\listoffigures
\listoftables

%----------------------------------------------------------------------------------------
%	CUERPO DE TEXTO
%----------------------------------------------------------------------------------------

\newpage
\pagestyle{cuerpo}
\thispagestyle{cuerpo}

\section{Introducción}

Para saltar de línea hay  que escribir la barra hacia la izquierda. Ahora si \\
Para saltar a párrafo hay que insertar una línea en blanco\\

Para saltar a párrafo sin sangría hay que insertar la barra a la izquierda en dos lineas\\
\\
Y por aquí sigue la cosa.

\subsection{Otro apartado}

Notas para mi, hay que decir como elaborar las tablas y figuras, un poco del modo matemático y después como manejar la bibliografía.

\subsection{Más texto suelto}
	\lipsum 

%----------------------------------------------------------------------------------------
%	BIBLIOGRAFIA
%----------------------------------------------------------------------------------------

\bibliography{bibo.bib} %Nombre del fichero de bibliografia.bib
\bibliographystyle{apalike} %Estilo de la bibliografia

\end{document}
